\section{Result-Pair Comparability, Nine-Item Split Audit}\label{app:result-audit}
The complete audit covers 39 records from 26 families, with 33 mitigation records from 20 families and six contextual records. Its 741 unordered pairs include 133 related-cohort comparisons and 608 outside the declared cohorts. The 15 same-family pairs are nested within the 133.
\begin{table}[t]\centering\footnotesize
\caption{Complete record-pair descriptive comparability.}\label{tab:pair-descriptive}
\begin{tabularx}{\columnwidth}{@{}Yrrr@{}}\toprule
Subset & Compatible & Incompatible & Unresolved\\\midrule
Same family ($n=15$) &7 (46.7\%)&8 (53.3\%)&0\\
Related cohort ($n=133$) &17 (12.8\%)&105 (78.9\%)&11 (8.3\%)\\\bottomrule
\end{tabularx}\end{table}
Eight same-family pairs change the primitive, target or cost object. Of seven aligned-outcome pairs, four change the collection/execution population and three lack a verified paired-test binding. The artifact preserves the original 528 judgments and identifies all 31 added related-pair judgments. Cohort membership, source locators and individual reasons are explicit; outside-cohort pairs are excluded rather than treated as adverse evidence. Section~\ref{app:record-joins} reports all 351 record-local audit cells and positive joint-binding identities.

\subsection{Nine-item split reporting audit}
Table~\ref{tab:six-check-audit} uses R4 for defended fitting, R5 for a matched adaptation contrast, R7 for inference, R8 for utility and R9 for cost across both the 16 original evaluated-mitigation families (144 cells) and the 20 total evaluated-mitigation families (including \texttt{wu2026cacheprune}, \texttt{lee2026otro}, \texttt{chen2026sparleak} and \texttt{hadad2026shape}; 180 cells).
\begin{table}[t]\centering\footnotesize
\caption{Nine-item split reporting audit over the 16 original evaluated-mitigation families (and 20 total including the four October~21 mitigation additions in parentheses). R1--R9 describe evidence for scoped reuse claims; they are separate from deployment Checks~1--3.}\label{tab:six-check-audit}
\setlength{\tabcolsep}{3.5pt}
\begin{tabularx}{\columnwidth}{@{}Yrrrrr@{}}\toprule
Audit item & Bound & Partial & \NRR{} & \NRV{} & \NA{}\\\midrule
R1. Target \& protection polarity & 16 (20) & 0 (0) & 0 (0) & 0 (0) & 0 (0)\\
R2. Receiver fields \& linkage & 13 (17) & 3 (3) & 0 (0) & 0 (0) & 0 (0)\\
R3. Placement \& action dispatch & 14 (17) & 2 (3) & 0 (0) & 0 (0) & 0 (0)\\
R4. Protected-condition fit & 4 (4) & 4 (5) & 1 (2) & 0 (0) & 7 (9)\\
R5. Matched adapted test & 2 (2) & 2 (2) & 5 (7) & 0 (0) & 7 (9)\\
R6. Workloads \& generator shifts & 14 (18) & 2 (2) & 0 (0) & 0 (0) & 0 (0)\\
R7. Defended inference effect & 14 (17) & 0 (1) & 2 (2) & 0 (0) & 0 (0)\\
R8. Task or component utility & 2 (4) & 2 (2) & 10 (11) & 0 (0) & 2 (3)\\
R9. Mechanism or protocol cost & 10 (11) & 1 (3) & 4 (5) & 0 (0) & 1 (1)\\\bottomrule
\end{tabularx}
\par\smallskip\raggedright\footnotesize
Bound: explicitly established for the scoped claim. Partial: a supported element coexists with a proxy or unlinked operating point (e.g., AHB refit code without verified Table~1 run binding; WRA Table~11 visibility utility separate from Tables~8--9 decoys). Only CachePrune across these 20 families jointly binds R7, R8 and R9 in one configuration package. Current totals include SparLeak's matched component utility and the 11-cell reinspection in Section~\ref{app:source-reconciliation}; historical coding columns remain unchanged.
\end{table}

% Dangol R5 is capture-scoped; live refit remains NR-V. See main text.


The current \path{combined_result_records.json} makes the six expansion records explicit alongside the 33 preserved baseline records: 39 records from 26 distinct families, while the reporting audit covers 20 mitigation families.

\subsection{Current source-follow-up corrections}
The first source-follow-up pass resolved three of 14 unverified family cells. For the cloud-traffic benchmark, \S V-C1/Fig.~9(a) locates execution-trigger instructions before and after the original prompt and reports the resulting reformulation experiment \cite{pouryousef2026traffic}, establishing R3 placement. For Shape and Substance, \S IX-B1 derives padding cost from unprotected timing measurements without reporting protected fitting or a matched adaptation contrast \cite{hadad2026shape}; R4/R5 therefore become \NRR{}. The targeted reinspection in Section~\ref{app:source-reconciliation} resolves the remaining 11 family cells as four P and seven U, with no new B. These judgments distinguish reported profiling and defended outcomes from an explicitly linked protected fit or matched adaptation contrast. The current 180-cell family matrix thus has no question marks. Six record-specific bindings remain unresolved because a family-level judgment cannot be inherited by every execution record. Historical coding columns are preserved; versioned source locators and the correction overlay generate the current tables.

SparLeak's \S7.5/Table~4 pairs reconstruction success and ROUGE at the same model and perturbation intensity \cite{chen2026sparleak}. R8 accepts component as well as task utility, so the current overlay corrects its partial code to B. Its reported spy overhead is attack cost and is not promoted to defense cost. The R8 distribution becomes 4 B, 2 P, 11 U and 3 inapplicable; historical records remain unchanged.
